% EDITORIAL QUESTION: How does fractionalization of spin-1/2 into four Majorana partons resolve quantum frustration while preserving physical Hilbert space dimension?
\section{Introduction and Physical Fractionalization}
\label{sec:intro}

Frustrated quantum magnets with bond-directional exchange interactions, epitomized by the Kitaev honeycomb model \cite{Kitaev2006}, host topological quantum spin liquid (QSL) phases characterized by emergent fractionalized Majorana fermions and $\mathbb{Z}_2$ gauge fluxes (visons). In candidate materials such as $\alpha\text{-RuCl}_3$, $\mathrm{Na}_2\mathrm{IrO}_3$, and $\alpha\text{-Li}_2\mathrm{IrO}_3$, external Zeeman magnetic fields $\mathbf{h}$ and residual Heisenberg ($J$) and symmetric off-diagonal ($\Gamma, \Gamma'$) exchange interactions induce strong non-integrable competition, driving transitions between topological spin liquids and field-polarized states.

The central computational challenge in quantitative mean-field treatments of these systems is twofold:
\begin{enumerate}
\item \textbf{Hilbert Space Doubling and Gauge Redundancy:} Representing physical spin-$1/2$ operators via four Majorana fermions per site doubles the local state space from $\dim \mathcal{H}_{\mathrm{phys}, i} = 2$ to $\dim \mathcal{F}_i = 4$, introducing an unphysical gauge-redundant sector that must be rigorously projected out via local gauge constraints.
\item \textbf{Non-Linear Fixed-Point Multistability:} In the presence of a magnetic field, the self-consistent Hartree-Fock-Bogoliubov map $\Phi: \mathcal{X} \to \mathcal{X}$ defines a non-linear dynamical system on the orthogonal Grassmannian manifold bundle over the Brillouin zone. Conventional fixed-point Picard iterations suffer from severe divergence and critical slowing down, while complex multi-basin landscapes exhibit competitive first-order hysteresis and metastable saddles.
\end{enumerate}

In this paper, we resolve these fundamental challenges through four key contributions:
\begin{enumerate}
\item \textbf{Exact Algebraic Invariants and Flat-Band Gauge Neutrality (Section~\ref{sec:algebra_and_invariants}):} We establish the local $\mathfrak{so}(4) \cong \mathfrak{su}(2)_{\mathrm{spin}} \oplus \mathfrak{su}(2)_{\mathrm{gauge}}$ Lie algebra homomorphism and prove the Flat-Band Gauge Neutrality Theorem: in the zero-field Kitaev limit, the self-consistent gauge Lagrange multipliers vanish identically ($\boldsymbol{\lambda}^* \equiv \mathbf{0}$) due to destructive interference among flat localized bands. Furthermore, we characterize the Pearson walk quadrature scaling ($\mathcal{O}(N_k^{-1/2})$) governing numerical Brillouin zone discretization noise.
\item \textbf{Field Response, Empirical Weak-Field Regressions, and KKT Counter-Torque (Section~\ref{sec:field_response}):} Under magnetic fields along $[111]$, we locate a sharp first-order transition between the chiral QSL and the field-polarized phase at $h_c^* / |K| = 0.3325 \pm 0.0005$ with an energy hysteresis barrier of $\Delta E \approx 0.057 |K|$. In the weak-field regime, $C_3$ and time-reversal selection rules mandate cubic gap scaling $\Delta_K \propto h^3$ (with finite-mesh empirical fit $\Delta_K^{\mathrm{MF}} = (2.592 \pm 0.003) h^3/|K|^2$), while the gauge Lagrange multiplier follows the finite-mesh empirical linear relation $\|\boldsymbol{\lambda}^*\| \approx \sqrt{5} h$; closed-form analytical evaluation of both prefactors (\texttt{IDEA-001}) and full-$Q$ Feshbach--Schur resolvent comparisons (\texttt{K-019-A}) remain explicitly delimited open questions.
\item \textbf{Grassmannian Solvers and Statistical Pareto Frontiers (Section~\ref{sec:solvers_and_real_space}):} By benchmarking an ensemble of $N_{\mathrm{trials}} = 30$ independent $C_3$-symmetric random seed perturbations per solver with 95\% Wilson score confidence intervals, we map the Pareto efficiency frontier of non-linear solvers and supplement dual updates with SPD-floored Lehmann Hessian preconditioning. Full Broyden and Joint DIIS achieve robust basin survival ($R_{90} = 0.35$), outperforming Picard by $>35\times$ in basin radius and $4.26\times$ in convergence speed.
\item \textbf{Real-Space Scaling, Rational Matrix Sign Purification, and Pfaffian Fluxes (Section~\ref{sec:solvers_and_real_space}):} By reformulating pure-\texttt{scMFT} Majorana Hamiltonians on real-space torus clusters ($M \in [128, 1568]$), we eliminate complex arithmetic, halving matrix memory ($2.0\times$). Combining cubic Newton-Schulz and quintic Kenney--Laub Pad\'e $[3/2]$ matrix sign purification $\Omega = \ii \operatorname{sgn}(\ii \mathbf{A})$ with canonical real skew-Schur decomposition and Parlett--Reid 12-Majorana plaquette Pfaffians yields compute-bound BLAS-3 throughput ($114.64\text{ GFLOP/s}$) and robust vortex/disorder metrology at machine-precision accuracy ($4 \times 10^{-15}$).
\end{enumerate}
